\section{Le cadre de travail Scrum}

\begin{frame}{Philosophie générale de Scrum}
    \begin{columns}
        \begin{column}{.55\textwidth}
            \begin{itemize}
                \item Scrum (« mêlée » en français) s'inspire des valeurs collectives du rugby.
                \item Approche incrémentale et itérative.
                \item Travail en équipe au cœur du projet.
            \end{itemize}
        \end{column}
        \begin{column}{.45\textwidth}
            \includegraphics[width=\linewidth]{rugby}
            \vfill
            \hfill{\scriptsize \textbf{Source :} Kiwi Flickr (CC BY 2.0)}\\[-.3em]
            \hfill{\scriptsize\url{https://flic.kr/p/3bnXZ}}
        \end{column}
    \end{columns}
\end{frame}

\definecolor{scrumrelease}{HTML}{005480}
\definecolor{scrumsprint}{HTML}{800014}
\definecolor{scrummilestone}{HTML}{806C00}

\tikzset{
    baseblock/.style={
        minimum height=3ex,
        rounded corners,
        fill=fg,
        text=white,
        inner sep=1.2ex,
    },
    block/.style={
        baseblock,
        text depth=0ex,
        text height=1.2ex,
    },
    release/.style={
        block,
        fill=scrumrelease,
    },
    sprint/.style={
        block,
        fill=scrumsprint,
    },
    point/.style={
        fill=fg,
        circle,
        inner sep=3pt,
    }
}

\begin{frame}[t]{Cadre temporel de Scrum}{Sprint : unité de base}
    \vspace{1.5cm}
    \centering\tikz\node[sprint] {Sprint};

    \begin{itemize}
        \item \textbf{Objectif : apporter un incrément.}
        \item Durée fixée à l’avance en semaines (2 -- 3 semaines).
        \item Échéance non modifiable : pas de retard, budget fixé.
        \item L’équipe reste la même.
    \end{itemize}
\end{frame}

\begin{frame}[t]{Cadre temporel de Scrum}{Déroulement d'un sprint}
    \vspace{1.5cm}
    \begin{tikzpicture}
        \fill[color=scrumsprint!50]
            (.5cm, 0) {[rounded corners] --
            (.5cm, 2.15) --
            (\linewidth - 2pt - 1cm, 2.15)} --
            (\linewidth - 2pt - 1cm, 0);
        \draw[arrow, very thick, ->] (0, 0) -- (\linewidth - 2pt, 0);

        \node[sprint] (sprint-master) at (.5\linewidth - .5pt, 2.15) {Sprint};

        \node[block, anchor=west] (sprint-plan) at (.5, 0) {Planification};
        \node[block, anchor=east, align=center, visible on=<4->] (sprint-retro) at (\linewidth - 2pt - 1cm, 0) {Rétrospective};
        \node[baseblock, left=.6cm of sprint-retro, align=center, visible on=<3->] (sprint-revue) {Revue de\\sprint};
        \node[baseblock, align=center, right=.6cm of sprint-plan, visible on=<2->] (sprint-daily) {Mêlées\\quotidiennes};

        \draw[arrow, very thick, ->, visible on=<2->] (sprint-daily) edge[in=-110,out=-70,loop,looseness=4] (sprint-daily);
    \end{tikzpicture}
\end{frame}

\begin{frame}[t]{Cadre temporel de Scrum}{Release : agrégat de sprints}
    \vspace{1.5cm}
    \centering\tikz\node[release] {Release};

    \begin{itemize}
        \item \textbf{Objectif : livrer une version.}
        \item Durée flexible (3 -- 4 mois).
    \end{itemize}
\end{frame}

\begin{frame}[t]{Cadre temporel de Scrum}{Déroulement d'une release}
    \vspace{1.5cm}
    \begin{tikzpicture}
        \fill[color=scrumrelease!50]
            (.5cm, 0) {[rounded corners] --
            (.5cm, 2) --
            (\linewidth - 2pt - 1cm, 2)} --
            (\linewidth - 2pt - 1cm, 0);
        \draw[arrow, very thick, ->] (0, 0) -- (\linewidth - 2pt, 0);

        \node[release] (release-master) at (.5\linewidth - .5pt, 2) {Release};

        \node[block, anchor=west] (release-prep) at (.5, 0) {Préparation};
        \node[sprint, visible on=<2->] (release-sprint) at (.5\linewidth, 0) {Sprints};
        \node[block, anchor=east, visible on=<3->] (release-final) at (\linewidth - 2pt - 1cm, 0) {Finalisation};

        \draw[arrow, very thick, ->, visible on=<2->] (release-sprint) edge[in=-120,out=-60,loop,looseness=5] (release-sprint);
    \end{tikzpicture}
\end{frame}

\begin{frame}{Cadre temporel de Scrum}{Déroulement d'un projet}
    \pgfdeclarelayer{background}
    \pgfdeclarelayer{front}
    \pgfsetlayers{background,front}

    \begin{tikzpicture}
        \begin{pgfonlayer}{front}
            \node[point] (start) {};
            \node[below=0cm of start, align=center] {Début du\\projet};

            \node[release, right=1cm of start] (release1) {\color{white}{Release 1}};
            \node[release, right=1cm of release1] (release2) {\color{white}{Release 2}};
            \node[release, right=1cm of release2] (release3) {\color{white}{Release 3}};

            \draw[
                decoration={brace, raise=.5cm},
                thick, decorate
            ] (release2.east) --
                node[below=.75cm] {3 – 4 mois}
                (release2.west);

            \node[point] (end) at (\linewidth - 2cm, 0) {};
            \node[below=0cm of end, align=center] {Fin du\\projet};
        \end{pgfonlayer}

        \begin{pgfonlayer}{background}
            \draw[very thick] (start.center) -- (release3.west);
            \draw[arrow, very thick, ->, dash pattern=on 6pt off 6pt]
                (release3.west) -- (end);
        \end{pgfonlayer}
    \end{tikzpicture}
    \vfill
    \begin{itemize}
        \item À la différence des sprints, les releases sont espacées.
        \item Permet de faire une rétrospective et de faire la fête.
    \end{itemize}
\end{frame}

\begin{frame}{Cadre temporel de Scrum}{En bref}
    \begin{itemize}
        \item Projet découpé en \tikz[baseline=(rel.base)] \node[release, scale=.7] (rel) {Releases}; $\rightarrow$ versions livrables.
        \item Releases découpées en \tikz[baseline=(spr.base)] \node[sprint, scale=.7] (spr) {Sprints}; $\rightarrow$ incréments au produit.
        \item Réunions régulières pour affiner le cycle de développement du produit.
        \item Coordination avec le client après chaque sprint.
    \end{itemize}
\end{frame}

\begin{frame}{Cadre humain de Scrum}{L'équipe selon Scrum}
    \textbf{Fonctionnement}

    \begin{itemize}
        \item L'équipe a le pouvoir sur la réalisation du projet.
        \item Équipe non-hiérarchisée, auto-organisée et indépendante.
    \end{itemize}

    \pause
    \vfill
    \textbf{Critères de choix de l'équipe}

    \begin{itemize}
        \item Taille de l'équipe : 2 à 9 personnes (optimum à 7).
        \item Équipe pluridisciplinaire.
        \item Cohésion et proximité géographique.
    \end{itemize}
\end{frame}

\begin{frame}{Cadre humain de Scrum}{Vue d'ensemble des rôles}
    % Schéma des différents rôles occupés par les acteurs de Scrum

\begin{tikzpicture}
    % L'équipe : ScrumMaster et programmeurs
    \foreach \prognum in {1,2,3,4,5}{
        \node[image, inner sep=0, outer sep=0] (prog-\prognum) at (\prognum*360/5+18 : 2.4cm and 1.8cm)
            {\includegraphics[width=1cm]{icon-coding}};
    }

    \draw ($(prog-3.south)!0.5!(prog-4.south)$)
        node[below] (prog-text) {Programmeurs};

    \node[image] (sm) {\includegraphics[width=1cm]{icon-trainer}\\ScrumMaster};

    \node[
        draw=fg, thick, ellipse,
        inner xsep=0cm, inner ysep=-.6cm,
        fit=(prog-1) (prog-2) (prog-3) (prog-4) (prog-5) (prog-text) (sm)
    ] (team) {}
        node[above] at (team.north) {\scshape\bfserieséquipe};

    % Interface du PO
    \node[
        rectangle, fill=white,
        yshift=-.25cm, inner xsep=1cm, inner ysep=1.5cm
    ] at (team.east) {};

    \node[image] (po) at (sm-|team.east)
        {\includegraphics[width=1cm]{icon-user}}
        node[anchor=north, align=center] at (po.south) {Product\\Owner};

    % Parties prenantes
    \node[image, right=of po] (pp2)
        {\includegraphics[width=1cm]{icon-woman}};

    \node[image, above=0cm of pp2] (pp1)
        {\includegraphics[width=1cm]{icon-man}};

    \node[image, below=0cm of pp2] (pp3)
        {\includegraphics[width=1cm]{icon-man}\\Parties\\prenantes};

    % Experts
    \node [image, left=of sm-|team.west] (exp2)
        {\includegraphics[width=1cm]{icon-scientist}};

    \node [image, above=0cm of exp2] (exp1)
        {\includegraphics[width=1cm]{icon-scientist}};

    \node [image, below=0cm of exp2] (exp3)
        {\includegraphics[width=1cm]{icon-scientist}\\Experts};
\end{tikzpicture}

\end{frame}

\begin{frame}{Cadre humain de Scrum}
\begin{columns}
    \begin{column}{.4\textwidth}
        \hspace{2em}\textbf{Scrum Master}
        \begin{itemize}
            \item Coach de l'équipe.
            \item Au service de l'équipe\\(pas un supérieur).
            \item Garant de l'application\\de Scrum.
            \item Aide à éliminer les obstacles.
        \end{itemize}
    \end{column}%
    \begin{column}{.6\textwidth}
        \newcommand{\scrumfocus}{sm}
        % Schéma des différents rôles occupés par les acteurs de Scrum,
% sans les labels pour prendre moins de place

\newcommand{\opacityat}[2]{
    \ifthenelse{\equal{\scrumfocus}{#1}}%
        {\begin{scope}[opacity=1, transparency group]#2\end{scope}}%
        {\begin{scope}[opacity=.3, transparency group]#2\end{scope}}
}

\begin{tikzpicture}
    % L'équipe : ScrumMaster et programmeurs
    \foreach \prognum in {1,2,3,4,5}{
        \opacityat{prog}{
            \node[image, inner sep=0, outer sep=0]
                (prog-\prognum) at (\prognum*360/5+18 : 1.4cm and 1.2cm)
                {\includegraphics[width=.8cm]{icon-coding}};
        }
    }

    \opacityat{sm}{
        \node[image] (sm) {\includegraphics[width=.8cm]{icon-trainer}};
    }

    \node[
        draw=fg, thick, ellipse,
        inner sep=0cm,
        fit=(prog-1) (prog-2) (prog-3) (prog-4) (prog-5) (sm)
    ] (team) {};

    % Interface du PO
    \node[
        rectangle, fill=white,
        inner xsep=.6cm, inner ysep=.6cm
    ] at (sm-|team.east) {};

    \opacityat{po}{
        \node[image] (po) at (sm-|team.east)
            {\includegraphics[width=.8cm]{icon-user}};
    }

    % Parties prenantes
    \opacityat{pp}{
        \node[image, right=.2cm of po] (pp2)
            {\includegraphics[width=.8cm]{icon-woman}};

        \node[image, above=0cm of pp2] (pp1)
            {\includegraphics[width=.8cm]{icon-man}};

        \node[image, below=0cm of pp2] (pp3)
            {\includegraphics[width=.8cm]{icon-man}};
    }

    % Experts
    \opacityat{exp}{
        \node [image, left=.2cm of sm-|team.west] (exp2)
            {\includegraphics[width=.8cm]{icon-scientist}};

        \node [image, above=0cm of exp2] (exp1)
            {\includegraphics[width=.8cm]{icon-scientist}};

        \node [image, below=0cm of exp2] (exp3)
            {\includegraphics[width=.8cm]{icon-scientist}};
    }
\end{tikzpicture}

    \end{column}
\end{columns}
\end{frame}

\begin{frame}{Cadre humain de Scrum}
\begin{columns}
    \begin{column}{.4\textwidth}
        \hspace{2em}\textbf{Product Owner}
        \begin{itemize}
            \item Médiateur entre l'équipe et les parties prenantes.
            \item Vision globale du produit (fonctionnalités, attentes, objectifs).
            \item Gère et priorise les tâches et répartit le temps.
        \end{itemize}
    \end{column}%
    \begin{column}{.6\textwidth}
        \newcommand{\scrumfocus}{po}
        % Schéma des différents rôles occupés par les acteurs de Scrum,
% sans les labels pour prendre moins de place

\newcommand{\opacityat}[2]{
    \ifthenelse{\equal{\scrumfocus}{#1}}%
        {\begin{scope}[opacity=1, transparency group]#2\end{scope}}%
        {\begin{scope}[opacity=.3, transparency group]#2\end{scope}}
}

\begin{tikzpicture}
    % L'équipe : ScrumMaster et programmeurs
    \foreach \prognum in {1,2,3,4,5}{
        \opacityat{prog}{
            \node[image, inner sep=0, outer sep=0]
                (prog-\prognum) at (\prognum*360/5+18 : 1.4cm and 1.2cm)
                {\includegraphics[width=.8cm]{icon-coding}};
        }
    }

    \opacityat{sm}{
        \node[image] (sm) {\includegraphics[width=.8cm]{icon-trainer}};
    }

    \node[
        draw=fg, thick, ellipse,
        inner sep=0cm,
        fit=(prog-1) (prog-2) (prog-3) (prog-4) (prog-5) (sm)
    ] (team) {};

    % Interface du PO
    \node[
        rectangle, fill=white,
        inner xsep=.6cm, inner ysep=.6cm
    ] at (sm-|team.east) {};

    \opacityat{po}{
        \node[image] (po) at (sm-|team.east)
            {\includegraphics[width=.8cm]{icon-user}};
    }

    % Parties prenantes
    \opacityat{pp}{
        \node[image, right=.2cm of po] (pp2)
            {\includegraphics[width=.8cm]{icon-woman}};

        \node[image, above=0cm of pp2] (pp1)
            {\includegraphics[width=.8cm]{icon-man}};

        \node[image, below=0cm of pp2] (pp3)
            {\includegraphics[width=.8cm]{icon-man}};
    }

    % Experts
    \opacityat{exp}{
        \node [image, left=.2cm of sm-|team.west] (exp2)
            {\includegraphics[width=.8cm]{icon-scientist}};

        \node [image, above=0cm of exp2] (exp1)
            {\includegraphics[width=.8cm]{icon-scientist}};

        \node [image, below=0cm of exp2] (exp3)
            {\includegraphics[width=.8cm]{icon-scientist}};
    }
\end{tikzpicture}

    \end{column}
\end{columns}
\end{frame}

\begin{frame}{Cadre humain de Scrum}
\begin{columns}
    \begin{column}{.4\textwidth}
        \hspace{2em}\textbf{Programmeurs}
        \begin{itemize}
            \item Effectuent les tâches du projet.
            \item Partagent leurs compétences et les mettent au service de l'équipe.
            \item Participent à la prise de décisions.
        \end{itemize}
    \end{column}%
    \begin{column}{.6\textwidth}
        \newcommand{\scrumfocus}{prog}
        % Schéma des différents rôles occupés par les acteurs de Scrum,
% sans les labels pour prendre moins de place

\newcommand{\opacityat}[2]{
    \ifthenelse{\equal{\scrumfocus}{#1}}%
        {\begin{scope}[opacity=1, transparency group]#2\end{scope}}%
        {\begin{scope}[opacity=.3, transparency group]#2\end{scope}}
}

\begin{tikzpicture}
    % L'équipe : ScrumMaster et programmeurs
    \foreach \prognum in {1,2,3,4,5}{
        \opacityat{prog}{
            \node[image, inner sep=0, outer sep=0]
                (prog-\prognum) at (\prognum*360/5+18 : 1.4cm and 1.2cm)
                {\includegraphics[width=.8cm]{icon-coding}};
        }
    }

    \opacityat{sm}{
        \node[image] (sm) {\includegraphics[width=.8cm]{icon-trainer}};
    }

    \node[
        draw=fg, thick, ellipse,
        inner sep=0cm,
        fit=(prog-1) (prog-2) (prog-3) (prog-4) (prog-5) (sm)
    ] (team) {};

    % Interface du PO
    \node[
        rectangle, fill=white,
        inner xsep=.6cm, inner ysep=.6cm
    ] at (sm-|team.east) {};

    \opacityat{po}{
        \node[image] (po) at (sm-|team.east)
            {\includegraphics[width=.8cm]{icon-user}};
    }

    % Parties prenantes
    \opacityat{pp}{
        \node[image, right=.2cm of po] (pp2)
            {\includegraphics[width=.8cm]{icon-woman}};

        \node[image, above=0cm of pp2] (pp1)
            {\includegraphics[width=.8cm]{icon-man}};

        \node[image, below=0cm of pp2] (pp3)
            {\includegraphics[width=.8cm]{icon-man}};
    }

    % Experts
    \opacityat{exp}{
        \node [image, left=.2cm of sm-|team.west] (exp2)
            {\includegraphics[width=.8cm]{icon-scientist}};

        \node [image, above=0cm of exp2] (exp1)
            {\includegraphics[width=.8cm]{icon-scientist}};

        \node [image, below=0cm of exp2] (exp3)
            {\includegraphics[width=.8cm]{icon-scientist}};
    }
\end{tikzpicture}

    \end{column}
\end{columns}
\end{frame}

\begin{frame}{Cadre humain de Scrum}
\begin{columns}
    \begin{column}{.4\textwidth}
        \hspace{2em}\textbf{Parties prenantes}
        \begin{itemize}
            \item Personnes intéressées par les résultats (clients, utilisateurs, managers).
            \item Participent aux revues.
        \end{itemize}
    \end{column}%
    \begin{column}{.6\textwidth}
        \newcommand{\scrumfocus}{pp}
        % Schéma des différents rôles occupés par les acteurs de Scrum,
% sans les labels pour prendre moins de place

\newcommand{\opacityat}[2]{
    \ifthenelse{\equal{\scrumfocus}{#1}}%
        {\begin{scope}[opacity=1, transparency group]#2\end{scope}}%
        {\begin{scope}[opacity=.3, transparency group]#2\end{scope}}
}

\begin{tikzpicture}
    % L'équipe : ScrumMaster et programmeurs
    \foreach \prognum in {1,2,3,4,5}{
        \opacityat{prog}{
            \node[image, inner sep=0, outer sep=0]
                (prog-\prognum) at (\prognum*360/5+18 : 1.4cm and 1.2cm)
                {\includegraphics[width=.8cm]{icon-coding}};
        }
    }

    \opacityat{sm}{
        \node[image] (sm) {\includegraphics[width=.8cm]{icon-trainer}};
    }

    \node[
        draw=fg, thick, ellipse,
        inner sep=0cm,
        fit=(prog-1) (prog-2) (prog-3) (prog-4) (prog-5) (sm)
    ] (team) {};

    % Interface du PO
    \node[
        rectangle, fill=white,
        inner xsep=.6cm, inner ysep=.6cm
    ] at (sm-|team.east) {};

    \opacityat{po}{
        \node[image] (po) at (sm-|team.east)
            {\includegraphics[width=.8cm]{icon-user}};
    }

    % Parties prenantes
    \opacityat{pp}{
        \node[image, right=.2cm of po] (pp2)
            {\includegraphics[width=.8cm]{icon-woman}};

        \node[image, above=0cm of pp2] (pp1)
            {\includegraphics[width=.8cm]{icon-man}};

        \node[image, below=0cm of pp2] (pp3)
            {\includegraphics[width=.8cm]{icon-man}};
    }

    % Experts
    \opacityat{exp}{
        \node [image, left=.2cm of sm-|team.west] (exp2)
            {\includegraphics[width=.8cm]{icon-scientist}};

        \node [image, above=0cm of exp2] (exp1)
            {\includegraphics[width=.8cm]{icon-scientist}};

        \node [image, below=0cm of exp2] (exp3)
            {\includegraphics[width=.8cm]{icon-scientist}};
    }
\end{tikzpicture}

    \end{column}
\end{columns}
\end{frame}

\begin{frame}{Cadre humain de Scrum}
\begin{columns}
    \begin{column}{.4\textwidth}
        \hspace{2em}\textbf{Experts}
        \begin{itemize}
            \item Jokers utilisés ponctuellement.
            \item Compétents dans un domaine particulier.
            \item Permettent d'outrepasser des difficultés.
        \end{itemize}
    \end{column}%
    \begin{column}{.6\textwidth}
        \newcommand{\scrumfocus}{exp}
        % Schéma des différents rôles occupés par les acteurs de Scrum,
% sans les labels pour prendre moins de place

\newcommand{\opacityat}[2]{
    \ifthenelse{\equal{\scrumfocus}{#1}}%
        {\begin{scope}[opacity=1, transparency group]#2\end{scope}}%
        {\begin{scope}[opacity=.3, transparency group]#2\end{scope}}
}

\begin{tikzpicture}
    % L'équipe : ScrumMaster et programmeurs
    \foreach \prognum in {1,2,3,4,5}{
        \opacityat{prog}{
            \node[image, inner sep=0, outer sep=0]
                (prog-\prognum) at (\prognum*360/5+18 : 1.4cm and 1.2cm)
                {\includegraphics[width=.8cm]{icon-coding}};
        }
    }

    \opacityat{sm}{
        \node[image] (sm) {\includegraphics[width=.8cm]{icon-trainer}};
    }

    \node[
        draw=fg, thick, ellipse,
        inner sep=0cm,
        fit=(prog-1) (prog-2) (prog-3) (prog-4) (prog-5) (sm)
    ] (team) {};

    % Interface du PO
    \node[
        rectangle, fill=white,
        inner xsep=.6cm, inner ysep=.6cm
    ] at (sm-|team.east) {};

    \opacityat{po}{
        \node[image] (po) at (sm-|team.east)
            {\includegraphics[width=.8cm]{icon-user}};
    }

    % Parties prenantes
    \opacityat{pp}{
        \node[image, right=.2cm of po] (pp2)
            {\includegraphics[width=.8cm]{icon-woman}};

        \node[image, above=0cm of pp2] (pp1)
            {\includegraphics[width=.8cm]{icon-man}};

        \node[image, below=0cm of pp2] (pp3)
            {\includegraphics[width=.8cm]{icon-man}};
    }

    % Experts
    \opacityat{exp}{
        \node [image, left=.2cm of sm-|team.west] (exp2)
            {\includegraphics[width=.8cm]{icon-scientist}};

        \node [image, above=0cm of exp2] (exp1)
            {\includegraphics[width=.8cm]{icon-scientist}};

        \node [image, below=0cm of exp2] (exp3)
            {\includegraphics[width=.8cm]{icon-scientist}};
    }
\end{tikzpicture}

    \end{column}
\end{columns}
\end{frame}

\begin{frame}{Gestion du travail avec Scrum}
    \begin{columns}
        \begin{column}{.69\textwidth}
            \begin{itemize}
                \item \textbf{Backlog :} outil de gestion dynamique des tâches.
                \item Représente le produit sous forme de tâches priorisées.
                \item Unique, public et visible par tous.
                \item Évolue au cours du temps.
            \end{itemize}
        \end{column}

        \begin{column}{.31\textwidth}
            \includegraphics[width=\linewidth]{backlog}

            \hfill{\scriptsize \textbf{Source :} Logan Ingalls (CC BY 2.0)}\\[-.3em]
            \hfill{\scriptsize\url{https://flic.kr/p/91QewP}}
        \end{column}
    \end{columns}
\end{frame}
